\documentclass[10pt,a4paper]{article}
\usepackage[utf8]{inputenc}
\usepackage[T1]{fontenc}
\usepackage[spanish,es-nodecimaldot]{babel}
\usepackage{lmodern}
\usepackage[margin=1.85cm,headheight=14pt]{geometry}
\usepackage{xcolor}
\usepackage{booktabs,tabularx,array,longtable}
\usepackage{enumitem}
\usepackage{fancyhdr}
\usepackage{hyperref}
\usepackage{listings}
\usepackage{microtype}
\definecolor{navy}{HTML}{173F67}
\definecolor{blue}{HTML}{276FA8}
\definecolor{soft}{HTML}{F3F6F9}
\definecolor{line}{HTML}{D5DFE8}
\hypersetup{colorlinks=true,linkcolor=navy,urlcolor=blue}
\pagestyle{fancy}
\fancyhf{}
\lhead{\small Seminario de Programación - Periodo 2}
\rhead{\small Instituto Jorge Robledo}
\cfoot{\small \thepage}
\setlength{\parindent}{0pt}
\setlength{\parskip}{5pt}
\setlist[itemize]{leftmargin=5mm,itemsep=2pt,topsep=3pt}
\setlist[enumerate]{leftmargin=6mm,itemsep=3pt,topsep=3pt}
\newcommand{\ruleline}{\par\noindent\color{line}\rule{\linewidth}{0.6pt}\color{black}\par}
\newcommand{\statusbox}[2]{\noindent\fcolorbox{line}{soft}{\parbox{0.94\linewidth}{\textbf{#1}\\#2}}}
\newcommand{\scorebox}[2]{\fcolorbox{line}{soft}{\parbox{0.43\linewidth}{\centering\textcolor{navy}{\Large\textbf{#1}}\\[-1mm]\small #2}}}
\lstset{basicstyle=\ttfamily\small,backgroundcolor=\color{soft},frame=single,rulecolor=\color{line},breaklines=true,columns=fullflexible,showstringspaces=false}
\begin{document}
\begin{center}
{\color{navy}\Large\bfseries Guía individual de avance}\\[2mm]
{\small Alejandro Rincón Torres - 11-C}
\end{center}
\ruleline

\begin{center}
\scorebox{No evaluable}{Resultado técnico provisional}\hfill
\scorebox{Sin porcentaje}{Avance frente a la rúbrica}
\end{center}
\statusbox{Proyecto principal}{Calculadora y panel de notas}
\statusbox{Repositorio y versión}{Pendiente de vincular - commit Sin SHA}
\statusbox{Nivel actual}{Sin evidencia verificable}

\section*{1. Lectura del estado actual}
La referencia actual es Carrito de restaurante de Jerónimo Rodríguez Peña, con 3.55/5.00 (71\%). La valoración automática analiza evidencia observable y debe confirmarse mediante ejecución, explicación y modificación en vivo.
\statusbox{Estado de evaluación}{No existe repositorio vinculado. No se asigna una nota técnica automática porque no hay código, commit ni ejecución verificables.}
\section*{2. Evidencias favorables}
\emph{Sin evidencia disponible.}
\section*{3. Riesgos o bloqueadores}
\begin{itemize}
\item No hay código verificable para analizar.
\end{itemize}


\newpage
\section*{4. Qué falta para alcanzar el nivel de referencia}
\begin{itemize}
\item Repositorio y commit inicial verificables.
\item Página HTML visible y organizada.
\item JavaScript con funciones, eventos y resultados en el DOM.
\item Validación, flujo completo, CSS, README y commits descriptivos.
\end{itemize}

\section*{5. Plan de trabajo para hoy}
\begin{enumerate}
\item Crear o vincular el repositorio del panel de notas.
\item Crear formulario de estudiante, materia y notas.
\item Validar escala, campos vacíos y valores no numéricos.
\item Calcular promedio y estado mediante funciones.
\item Renderizar resultados en una tabla y conservarlos con localStorage.
\item Agregar CSS, README y commits descriptivos.
\end{enumerate}
\section*{6. Distribución sugerida si trabaja en pareja}
\begin{itemize}
\item Integrante A: interfaz HTML y CSS.
\item Integrante B: JavaScript, DOM y persistencia.
\item Ambos: validaciones, pruebas, README y sustentación.
\end{itemize}

\statusbox{Regla de trabajo grupal}{El repositorio y el puntaje técnico pueden ser comunes, pero cada integrante debe aparecer en el README, tener un rol concreto, explicar su aporte y realizar un commit identificable o una modificación en vivo.}

\newpage
\section*{7. Evidencias que deben quedar subidas}
\begin{enumerate}
\item Código actualizado y ejecutable desde una página HTML principal.
\item README con propósito, integrantes, roles, instrucciones y estado actual.
\item Al menos tres commits descriptivos; evitar mensajes como ``final'', ``cambios'' o ``Add files via upload''.
\item Demostración del flujo principal y prueba de un caso de error.
\item Lista de lo terminado y de lo pendiente para la próxima clase.
\end{enumerate}
\section*{8. Verificación con el docente}
\begin{enumerate}
\item Abrir el commit actual y ejecutar la página principal.
\item Explicar entrada, procesamiento, estado y salida.
\item Identificar el bloqueador más importante.
\item Realizar una corrección pequeña en vivo.
\item Crear un commit descriptivo y volver a ejecutar.
\end{enumerate}
\section*{9. Ejemplos de commits válidos}
\begin{lstlisting}
feat: agrega formulario y resultados visibles
fix: valida cantidades y evita valores negativos
style: organiza la interfaz con grid responsive
docs: documenta integrantes, roles y ejecucion
\end{lstlisting}
\vfill
{\small\textit{Generado con el estado observado en la copia central. La valoración técnica no sustituye la sustentación individual.}}
\end{document}
